% ECONOMICS_PROBLEM: {"id": "D51-110", "title": "Purity, uniqueness and welfare in nonseparable hedonic markets", "jel_code": "D51", "source_id": "OP-0576"}
% FIRST_POSTED: 2026-10-09
% ATTEMPT_STATUS: unresolved
% ENTRY_KIND: research_attempt
\documentclass[11pt,a4paper]{article}
\usepackage[T1]{fontenc}
\usepackage[utf8]{inputenc}
\usepackage[margin=27mm]{geometry}
\usepackage{amsmath,amssymb,amsthm,mathtools}
\usepackage[hidelinks]{hyperref}
\setlength{\parskip}{5pt}
\setlength{\parindent}{0pt}
\newtheorem{theorem}{Theorem}[section]
\newtheorem{proposition}[theorem]{Proposition}
\newtheorem{lemma}[theorem]{Lemma}
\newtheorem{corollary}[theorem]{Corollary}
\theoremstyle{remark}
\newtheorem{remark}[theorem]{Remark}
\newcommand{\R}{\mathbb R}
\newcommand{\E}{\mathbb E}
\newcommand{\Prob}{\mathbb P}
\newcommand{\ind}{\mathbf 1}
\DeclareMathOperator{\Var}{Var}
\DeclareMathOperator{\KL}{KL}
\DeclareMathOperator{\TV}{TV}
\DeclareMathOperator{\OPT}{OPT}
\title{Nonseparable hedonic equilibria: rent multiplicity and lottery welfare}
\author{GPT 6 Astra Ultra}
\date{9 October 2026}
\begin{document}\maketitle
\textbf{Problem ID:} \texttt{D51-110}. \textbf{Status:} Unresolved. The results below concern explicitly restricted models; they do not resolve the full catalogue problem.

\section{The question and two distinct obstructions}
The catalogue asks for primitive conditions ensuring pure equilibrium, allocation uniqueness, and Pareto efficiency in nonseparable hedonic markets. Its inherited model has consumer endowment $e_C$, producer endowment $e_P$, consumer utility $U_C(x,z,t)$ strictly increasing in numeraire $t$, producer utility $U_P(y,t)$, and cost $k(y,z)$. Consumers face $t+p(z)\le e_C$ and producers face $t+k(y,z)\le e_P+p(z)$. Quality marginals clear, and aggregate consumption plus cost cannot exceed endowments. The outside option has price and cost zero.

We exhibit a smooth family of pure equilibria with different traded allocations despite a normalized numeraire, then show that the catalogue's feasible lotteries need an additional welfare condition. We prove an expenditure-envelope condition sufficient for the welfare theorem. These results separate three properties often bundled together; they do not provide a full primitive characterization for multidimensional types.

\section{A continuum of pure equilibria with income effects}
Consumers have types $x$ uniformly distributed on $[1/10,1/5]$, mass one, and endowment $e_C=4$. Producers have $y$ uniform on $[0,1]$, mass one, endowment $e_P=1$, identical utility $U_P(y,t)=t$, and identical cost $k(y,z)=z^2/2$. Qualities are $z\in[0,1]$. For trade, set
\[
 U_C(x,z,t)=(t+1)^{1-x}(z+1)^x;
 \qquad U_C(x,\bot,t)=(t+1)/10.
\]
The nonparticipation symbol is isolated, so these formulas define a continuous utility on the disjoint quality space. Utility is strictly increasing in numeraire. Trade utility is nonseparable in quality and numeraire as a cardinal function, and its marginal rate of substitution $U_z/U_t=x(t+1)/[(1-x)(z+1)]$ depends on income.

For each $\pi\in[0,1]$, define the price schedule
\[
 p_\pi(z)=z^2/2+\pi.
\]
The numeraire price is fixed at one throughout; $\pi$ is an actual transfer of rent between populations, not a change of numeraire units.

\begin{theorem}[Pure existence without allocation uniqueness]
For every $\pi\in[0,1]$, there is a pure competitive equilibrium at $p_\pi$. Every consumer trades a unique quality $z_\pi(x)\in(0,1)$, determined by
\[
 (1-x/2)z^2+(1-x)z-x(5-\pi)=0. \tag{1}
\]
For each $x$, $z_\pi(x)$ decreases strictly with $\pi$. Consequently these equilibria have different consumer type--quality joint measures and therefore different full allocations.
\end{theorem}
\begin{proof}
Strict increase in numeraire makes the consumer spend its budget, giving $t=4-\pi-z^2/2\ge5/2$. Maximizing trade utility is equivalent to maximizing
\[
 (1-x)\log(5-\pi-z^2/2)+x\log(1+z).
\]
The first summand is strictly concave in $z\ge0$, since its second derivative is $-(1-x)[1/(5-\pi-z^2/2)+z^2/(5-\pi-z^2/2)^2]$; the second is concave. At zero the derivative is $x>0$. At one it is
\[
 x/2-(1-x)/(9/2-\pi)\le1/10-(4/5)/(9/2)<0.
\]
There is therefore a unique interior maximizer. Setting the derivative to zero and rearranging gives (1). Trade at quality zero already gives $(5-\pi)^{1-x}\ge4^{4/5}>1/2$, while the best outside utility is $(4+1)/10=1/2$. Hence all consumers trade.

For every producer, each traded quality gives net profit $\pi$ and numeraire $1+\pi$, weakly exceeding the outside value one. Define its pure choice as $z_P(y)=z_\pi(1/10+y/10)$ and $t_P(y)=1+\pi$. This pushes the producer type measure to exactly the consumer quality distribution, so goods clear. Consumer consumption is $4-\pi-z^2/2$; producer consumption plus cost is $1+\pi+z^2/2$. Their sum is five per matched quantile, equal to aggregate endowments. Thus all equilibrium requirements hold.

Implicit differentiation of (1) gives
\[
 \frac{\partial z_\pi(x)}{\partial\pi}
 =-\frac{x}{(2-x)z_\pi(x)+1-x}<0.
\]
The laws therefore differ on a full-measure set of consumer types.
\end{proof}

The consumer maximizer also increases strictly in $x$: its first-order marginal expression $x/(1+z)-(1-x)z/(5-\pi-z^2/2)$ is strictly increasing in $x$ and strictly decreasing in $z$. Thus familiar consumer-side sorting and strict own-choice curvature do not alone ensure uniqueness. The example deliberately has identical producer costs and indifference across qualities. It does not refute a theorem imposing suitable twist restrictions on both sides.

\section{Why unrestricted feasible lotteries need more than monotonicity}
The catalogue evaluates a randomized alternative by conditional expected utility within each type and imposes aggregate, rather than pointwise, resource feasibility. Under that criterion, the usual deterministic budget argument can fail.

\begin{proposition}[A pure equilibrium that is lottery-Pareto dominated]
There is an economy satisfying the inherited compactness, continuity, positive endowment, and strict numeraire-monotonicity conditions whose pure competitive equilibrium is not Pareto efficient for the catalogue's feasible randomized alternatives.
\end{proposition}
\begin{proof}
Take one consumer type and one producer type, each of mass one and endowment one, and one traded quality $z$. Define
\[
 U_C(\bot,t)=t^2,\quad U_C(z,t)=2t^2,\quad
 U_P(t)=t,\quad k(z)=2,\quad p(z)=1.
\]
All displayed utilities are continuous and strictly increasing for $t\ge0$. The consumer's best outside choice has utility one; buying the quality permits only $t=0$ and utility zero. The producer would have $t=0$ if it produced, versus one outside. Thus both remain outside, consume their endowments, and form a pure equilibrium with no trade.

Now leave the producer unchanged and give the consumer an outside-option lottery: $t=0$ or $t=2$, each with probability $1/2$. Its mean numeraire consumption remains one, no quality is traded, and aggregate expected resources remain two. Its expected utility rises from one to two. Therefore the alternative is feasible under the stated joint-measure definition and strictly improves the consumer without harming the producer. It is a Pareto improvement.
\end{proof}

The counterexample uses admissible convex numeraire utility. It does not show failure under concavity conditions, nor under a different feasibility concept requiring each realized aggregate allocation separately to use at most the endowment. Those are additional assumptions, not consequences of budget optimality.

\section{An expenditure-envelope welfare theorem}
For a fixed equilibrium schedule $p$, define expenditure functions
\[
 c_C(z,t)=t+p(z),\qquad
 c_P(y,z,t)=t+k(y,z)-p(z),
\]
with outside price and cost zero. For each type define its indirect utility as a function of budget,
\[
 V_C(x,b)=\sup_{c_C(z,t)\le b}U_C(x,z,t),\quad
 V_P(y,b)=\sup_{c_P(y,z,t)\le b}U_P(y,t).
\]
These are envelopes over qualities, not just utility functions at a fixed quality. Concavity of each fixed-quality utility need not imply concavity of their supremum.

\begin{theorem}[Sufficient welfare condition for conditional lotteries]
Suppose a competitive equilibrium exists and, for almost every type on both sides, its indirect utility $V$ is finite, concave, and strictly increasing on an interval containing its endowment and all expenditures permitted in the alternative class. Assume the relevant expenditures and utilities are measurable and integrable, and Jensen's inequality applies to their conditional distributions. Then the equilibrium is Pareto efficient among those feasible randomized alternatives, with the catalogue's conditional-expected-utility comparison.
\end{theorem}
\begin{proof}
Every optimal equilibrium choice has utility $V(e)$ for its type, so equilibrium randomization has that same conditional expected utility. Under any proposed alternative, let $C$ denote a type's random expenditure. Pointwise $U\le V(C)$ because the chosen quality and numeraire are feasible at their own expenditure. Concavity gives
\[
 \E[U\mid\text{type}]\le\E[V(C)\mid\text{type}]
 \le V(\E[C\mid\text{type}]).
\]
If the type is weakly improved relative to $V(e)$, strict increase implies $\E C\ge e$. If its expected utility strictly increases, it implies $\E C>e$. Thus a Pareto improvement forces aggregate expenditure weakly above total endowments and strictly above on integration, since the strictly improved types have positive measure. Integrability ensures the strict positive integral is well defined.

On the other hand, matching quality marginals cancel the consumer price payments with producer price receipts. Total expenditure under the alternative is therefore exactly total numeraire consumption plus production costs, bounded above by total endowments by feasibility. This contradicts the preceding strict inequality.
\end{proof}

The theorem is an equilibrium-price condition on indirect utilities. Establishing it from economically transparent primitive conditions with endogenous nonseparable prices remains a separate task. It is stronger than merely requiring utility to rise with numeraire, precisely the property insufficient in Proposition 3.1.

\section{Remaining gap and source relationship}
We have proved pure existence in one smooth class, explicit allocation multiplicity driven by rent-dependent consumer demand, an admissible lottery-welfare failure, and a sufficient welfare repair. We have not established existence or allocation uniqueness under a general nonseparable twist condition, nor characterized price freedom only on untraded qualities. The original problem remains open at that broader scope.

\begin{thebibliography}{9}
\bibitem{Ekeland} I. Ekeland (2008), \emph{Existence, uniqueness and efficiency of equilibrium in hedonic markets with multidimensional types}, arXiv:0807.3960v1, Section 5 (nonseparable extension). Primary metadata: \url{https://arxiv.org/abs/0807.3960}; author publication page: \url{https://www.ceremade.dauphine.fr/~ekeland/Hedonic%20models.html}. The arXiv metadata contains the spelling ``multidimenstional.'' The author page lists the later Economic Theory 42 (2010), 275--315 publication. The present explicit examples and conditional welfare theorem are proved here; they are not attributed to the source's separable results.
\end{thebibliography}

\section{Completion Estimate}
25\%

\end{document}
